\chapter{Wie esim rechnet – und wo die Grenzen liegen}
\label{kap:numerik}

Für die Benutzung muss man die Numerik nicht kennen. Wer aber verstehen
will, warum eine Rechnung schnell oder langsam ist, oder Ergebnisse
einschätzen muss, findet hier das Wichtigste.

\section{Das Gleichungssystem}

Alles – Schaltung, eigene Gleichungen, Teilschaltungen – wird zu einem
einzigen System $\frac{d}{dt}q(x) + f(x,t) = 0$ zusammengesetzt. Jedes
Bauteil trägt seine Ströme zu $f$ und seine Ladungen und Flüsse zu $q$ bei,
dazu die Ableitungen $G = \partial f/\partial x$ und
$C = \partial q/\partial x$. Für eigene Gleichungen gehen
\code{der(...)}-Terme nach $q$, der Rest nach $f$; die Ableitungen kommen
aus automatischem Differenzieren.

\section{Arbeitspunkt}

Newton-Verfahren mit der Spannungsbegrenzung von SPICE an pn-Übergängen.
Scheitert es, folgen nacheinander:

\begin{enumerate}
  \item eine \textbf{Pseudo-Transiente}: Jeder Knoten ist gedämpft an seine
    vorherige Lösung gebunden, die Dämpfung wird schrittweise gelöst. Das
    folgt dem Einschaltvorgang und bewältigt auch lange Ketten von Stufen
    hoher Verstärkung;
  \item \textbf{gmin-Stepping}: zusätzliche Leitwerte nach Masse, die
    schrittweise verschwinden;
  \item \textbf{Quellen-Stepping}: alle Quellen werden von null aus
    hochgefahren.
\end{enumerate}

\section{Lineare Algebra}

Die Matrizen werden dünnbesetzt in LU-Faktoren zerlegt. Ihre Struktur wird
einmal analysiert; die Bauteile schreiben ihre Einträge danach über eine
gespeicherte Zuordnung direkt an ihren Platz. Ist die Matrix gegenüber dem
letzten Schritt unverändert (lineare Schaltung bei gleicher Schrittweite),
entfällt die Zerlegung ganz. Bei nichtlinearen Schaltungen wird die
zerlegte Matrix der ersten Newton-Iteration eines Zeitschritts
weiterverwendet, solange die Iteration gut konvergiert.

\section{Zeitintegration}

\begin{itemize}
  \item Implizite Integration; je Schritt ein Newton-Problem
    $f(x) + a\,q(x) + b = 0$.
  \item Die Schrittweite folgt aus dem Abbruchfehler der Ladungen $q$
    (Differenz zur Extrapolation ihres Verlaufs), nicht der Unbekannten
    selbst. So bleiben algebraische Größen außen vor, und schwach
    bestimmte, potentialfreie Schaltungsteile stören die
    Schrittweitensteuerung nicht.
  \item Mit \code{bdf} wird nach jedem Schritt geschätzt, welchen Fehler
    die nächstniedrigere und die nächsthöhere Ordnung gehabt hätten;
    gewählt wird die mit dem größten zulässigen Schritt. Die Schrittweite
    bleibt konstant, bis sie verdoppelt werden kann oder verkleinert werden
    muss (wie in DASSL).
  \item Knicke der Quellen, Ereignisse, Laufzeiten und Totzeiten sind
    Haltepunkte. Danach startet die Integration neu mit Ordnung~1; auch
    dieser erste Schritt ist fehlerkontrolliert.
\end{itemize}

\section{Ereignisse}

Jeder Vergleich in einem Ausdruck liefert eine Ereignisfunktion. Wechselt
eine davon in einem Schritt das Vorzeichen, wird der Schritt so lange
verkürzt (Sekantenverfahren mit Bisektion als Absicherung), bis er knapp
hinter dem Wechsel endet. Zwischen zwei Ereignissen sind die Ausgänge der
Vergleiche und die Stellungen der Schalter eingefroren – die Gleichungen
sind also glatt, auch für Newton. Danach folgen Zuweisungen und
Umschalten. Ein sehr kurzer Schritt, in dem alle Stellungen frei dem
Newton-Verfahren folgen, bringt voneinander abhängige Schaltelemente und die
algebraischen Unbekannten in Einklang; dann startet die Integration neu.

\section{AC und Rauschen}

Die AC-Analyse löst $(G + j\omega C)\,X = B$ mit den Ableitungen im
Arbeitspunkt. Für das Rauschen ist jede Rauschquelle ein Strom zwischen zwei
Knoten; ihre Übertragung zum Ausgang liefert eine adjungierte Lösung je
Frequenz, $(G + j\omega C)^T y = e_\text{aus}$. Die Beiträge werden
quadratisch summiert.

\section{Parallelität}

Studien verteilen ihre Läufe auf alle Kerne. Große Schaltungen (ab 1000
Bauteilen) werten zusätzlich ihre Bauteile parallel aus: Die Bauteile
werden in 16 Blöcke geteilt, die sich die Threads der Reihe nach holen;
jeder Block schreibt in einen eigenen Puffer, die Puffer werden in fester
Reihenfolge addiert. Deshalb sind die Ergebnisse für jede Zahl von Threads
bitgleich. Die Matrixzerlegung bleibt seriell und bestimmt dann die
Rechenzeit; gemessen sind rund 26\,\% Gewinn bei 20\,000 Dioden.

\section{Grenzen}

\begin{description}[style=nextline,leftmargin=2em]
  \item[Index]
    Die Verfahren lösen Systeme vom Index 1 und 2; höhere Indizes werden
    automatisch reduziert, soweit sich die Ableitungen als Ausdruck bilden
    lassen (Kapitel~\ref{kap:format-gleichungen}). Systeme vom Index 2 mit
    \code{"uic": true} brauchen Startwerte, die die Bedingungen erfüllen.
  \item[Ereignisse]
    Gleiten entlang einer Schaltfläche (Reibung als \code{if}, ein Schalter
    ohne Hysterese in einer Rückkopplung) führt zu endlosem Umschalten und
    wird abgebrochen.
  \item[Schalter und ideale Dioden]
    haben endliche Widerstände (\code{ron}, \code{roff}). Sehr große
    Verhältnisse \code{roff/ron} verschlechtern die Kondition.
  \item[Genauigkeit]
    \code{reltol} begrenzt den Fehler je Schritt, nicht den Gesamtfehler.
  \item[Schwach gedämpfte Schwingungen]
    BDF-Formeln der Ordnung 3 bis 5 sind dafür nicht stabil genug; die
    Ordnungssteuerung weicht dann auf niedrige, dämpfende Ordnungen aus.
    Für Schwingkreise hoher Güte \code{trap} verwenden.
  \item[Modelle]
    Klassische SPICE-Modelle, EKV und BSIM3v3 (ohne dessen Rauschmodelle
    und nicht-quasistatische Effekte). BSIM4 und neuere Modelle fehlen,
    ebenso der Skineffekt von Leitungen.
  \item[Große Schaltungen]
    Getestet bis 40\,000 Knoten. Ruhende Schaltungsteile werden nicht
    übersprungen; die Rechenzeit wächst mit der Zahl der Bauteile mal der
    Zahl der Zeitschritte.
  \item[Potentialfreie Schaltungsteile]
    brauchen einen Gleichstrompfad nach Masse. Ihre Spannungen gegen Masse
    sind dann nur auf Millivolt genau; Differenzen sind davon nicht
    betroffen.
\end{description}

\section{Erweitern}

Wer esim um ein Bauteil erweitern will, leitet in C++ von \code{Device} ab
und implementiert \code{stamp()}: Ströme, die einen Knoten verlassen, nach
$f$, Ladungen nach $q$, ihre Ableitungen nach $G$ und $C$. Zusätzliche
Unbekannte (innere Knoten, Zweigströme) werden in \code{setup()} angelegt;
eingetragen wird das Bauteil in \datei{src/io/Netlist.cpp}.
\datei{src/circuit/devices/Passive.cpp} und \datei{Diode.cpp} sind gute
Vorlagen. Für viele Fälle genügt allerdings eine Verhaltensquelle mit
Gleichungen in einer Teilschaltung.

Der numerische Kern ist auch ohne Eingabedatei nutzbar: von
\code{DaeSystem} ableiten, \code{size()} und \code{evaluate()}
implementieren und mit \code{runTransient()} lösen;
\datei{examples/dae/Robertson.hpp} zeigt das.
